%%%%%%%%%%%%%%%%%%%%%%%%%%%%%%%%%%%%%%%%%%%%%%%%%%%%%%%%%%%%%%%%%%%%%%%%%%%%%%%%
% DEPENDENCY ANALYSIS
%%%%%%%%%%%%%%%%%%%%%%%%%%%%%%%%%%%%%%%%%%%%%%%%%%%%%%%%%%%%%%%%%%%%%%%%%%%%%%%%

\chapter[Dependency Analysis]{\IconBox\ Dependency Analysis}

This chapter provides an overview of the software dependency ecosystem identified in the generated Software Bill of Materials (SBOM). The dependency analysis examines the composition of the \textbf{\ProjectName}'s software supply chain, including third-party libraries, package classifications, and dependency relationships within the agent codebase.

The generated SBOM was further analyzed using the \textbf{\VulScanner} Software Composition Analysis (SCA) tool to assess the security posture of these dependencies by identifying publicly disclosed vulnerabilities and available remediation information. Together, the SBOM and vulnerability assessment provide complete visibility into the software supply chain and its associated security risks.

\vspace{0.4cm}

%%%%%%%%%%%%%%%%%%%%%%%%%%%%%%%%%%%%%%%%%%%%%%%%%%%%%%%%%%%%%%%%%%%%%%%%%%%%%%%%
% DEPENDENCY OVERVIEW
%%%%%%%%%%%%%%%%%%%%%%%%%%%%%%%%%%%%%%%%%%%%%%%%%%%%%%%%%%%%%%%%%%%%%%%%%%%%%%%%

\begin{table}[H]
\centering
\small
\renewcommand{\arraystretch}{1.3}

\begin{tabular}{p{5cm}p{8cm}}
\toprule
\rowcolor{ZWPrimary}
\textbf{\color{white}Metric} & \textbf{\color{white}Value} \\
\midrule
Total Components & \TotalComponents \\
Library Components & \LibraryComponents \\
Application Components & \AppComponents \\
Package Manager & pip \\
Primary Ecosystem & Python \\
SBOM Format & \SbomFormat \\
Dependency Identification & Package URL (PURL) \\
\bottomrule
\end{tabular}
\caption{Dependency Overview}
\end{table}

\vspace{0.4cm}

%%%%%%%%%%%%%%%%%%%%%%%%%%%%%%%%%%%%%%%%%%%%%%%%%%%%%%%%%%%%%%%%%%%%%%%%%%%%%%%%
% DEPENDENCY COMPOSITION
%%%%%%%%%%%%%%%%%%%%%%%%%%%%%%%%%%%%%%%%%%%%%%%%%%%%%%%%%%%%%%%%%%%%%%%%%%%%%%%%

\section{Dependency Composition}

The generated SBOM indicates that the \ProjectName\ relies on third-party open-source libraries. These libraries provide functionality such as networking, system telemetry, cryptography, event logging, and process auditing.

Only a small number of components represent application artifacts, while the majority are reusable libraries obtained through the pip Python package index.

\vspace{0.4cm}

%%%%%%%%%%%%%%%%%%%%%%%%%%%%%%%%%%%%%%%%%%%%%%%%%%%%%%%%%%%%%%%%%%%%%%%%%%%%%%%%
% DEPENDENCY CATEGORIES
%%%%%%%%%%%%%%%%%%%%%%%%%%%%%%%%%%%%%%%%%%%%%%%%%%%%%%%%%%%%%%%%%%%%%%%%%%%%%%%%

\section{Major Dependency Categories}

\begin{table}[H]
\centering
\small
\renewcommand{\arraystretch}{1.3}

\begin{tabular}{p{5cm}p{8cm}}
\toprule
\rowcolor{ZWPrimary}
\textbf{\color{white}Category} & \textbf{\color{white}Examples} \\
\midrule
HTTP Communication & requests \\
Secure Sockets & python-socketio \\
Cryptography & cryptography \\
Compression & zstandard \\
System Telemetry & wmi, pywin32 \\
Compilation & nuitka \\
\bottomrule
\end{tabular}
\caption{Major Dependency Categories}
\end{table}

\vspace{0.4cm}

%%%%%%%%%%%%%%%%%%%%%%%%%%%%%%%%%%%%%%%%%%%%%%%%%%%%%%%%%%%%%%%%%%%%%%%%%%%%%%%%
% DEPENDENCY CHARACTERISTICS
%%%%%%%%%%%%%%%%%%%%%%%%%%%%%%%%%%%%%%%%%%%%%%%%%%%%%%%%%%%%%%%%%%%%%%%%%%%%%%%%

\section{Dependency Characteristics}

The dependency ecosystem exhibits the following characteristics:
\begin{itemize}
\item Large number of reusable open-source libraries.
\item Predominantly Python packages optimized for host operations.
\item Integration with the central administration platform APIs.
\item Extensive use of utility and OS-specific libraries.
\item Modern package management through pip.
\item Well-defined Package URLs (PURLs) for component identification.
\end{itemize}

\vspace{0.4cm}

%%%%%%%%%%%%%%%%%%%%%%%%%%%%%%%%%%%%%%%%%%%%%%%%%%%%%%%%%%%%%%%%%%%%%%%%%%%%%%%%
% SECURITY CONSIDERATIONS
%%%%%%%%%%%%%%%%%%%%%%%%%%%%%%%%%%%%%%%%%%%%%%%%%%%%%%%%%%%%%%%%%%%%%%%%%%%%%%%%

\section{Security Considerations}

Software dependencies represent one of the largest attack surfaces in modern applications. Every third-party library introduces potential security, maintenance, licensing, and operational risks that must be continuously monitored throughout the software development lifecycle.

The Software Bill of Materials (SBOM) generated using Syft provides a complete inventory of software components, while Grype and the ZeroWatch Agent Scan analyze those components against multiple vulnerability databases to identify known security issues.

Together, these tools enable organizations to:
\begin{itemize}
\item Identify vulnerable software packages.
\item Correlate dependencies with published CVEs and security advisories.
\item Detect components with available security updates.
\item Prioritize vulnerability remediation based on severity.
\item Assess software supply chain risk.
\item Monitor dependency changes across software releases.
\item Support continuous Software Composition Analysis (SCA).
\end{itemize}

\vspace{0.4cm}

%%%%%%%%%%%%%%%%%%%%%%%%%%%%%%%%%%%%%%%%%%%%%%%%%%%%%%%%%%%%%%%%%%%%%%%%%%%%%%%%
% BENEFITS
%%%%%%%%%%%%%%%%%%%%%%%%%%%%%%%%%%%%%%%%%%%%%%%%%%%%%%%%%%%%%%%%%%%%%%%%%%%%%%%%

\section{Benefits of Dependency Analysis}

Maintaining visibility into software dependencies provides several operational and security advantages:
\begin{itemize}
\item Improved software inventory management.
\item Complete visibility into direct and transitive dependencies.
\item Faster identification and remediation of vulnerable packages.
\item Better software governance and dependency lifecycle management.
\item Enhanced software supply chain transparency.
\item Simplified regulatory compliance and audit preparation.
\item Easier impact analysis during newly disclosed vulnerabilities.
\item Support for continuous Software Composition Analysis (SCA).
\item Improved release readiness through proactive dependency monitoring.
\end{itemize}

\vspace{0.4cm}

%%%%%%%%%%%%%%%%%%%%%%%%%%%%%%%%%%%%%%%%%%%%%%%%%%%%%%%%%%%%%%%%%%%%%%%%%%%%%%%%
% OBSERVATIONS
%%%%%%%%%%%%%%%%%%%%%%%%%%%%%%%%%%%%%%%%%%%%%%%%%%%%%%%%%%%%%%%%%%%%%%%%%%%%%%%%

\section{Key Observations}

The dependency analysis reveals that the \ProjectName\ relies extensively on reusable open-source components from the Python ecosystem. These dependencies provide essential functionality while significantly reducing development effort and accelerating feature delivery.

The use of standardized Package URLs (PURLs) and the CycloneDX specification improves component traceability and enables seamless integration with Software Composition Analysis (SCA) tools such as Grype and the ZeroWatch Agent Scan. This integration allows known security vulnerabilities to be correlated directly with software dependencies, supporting rapid vulnerability identification and remediation.

Continuous dependency monitoring, together with periodic SBOM generation and vulnerability assessment, forms an essential part of maintaining a secure and well-governed software supply chain.

\vspace{0.4cm}

%%%%%%%%%%%%%%%%%%%%%%%%%%%%%%%%%%%%%%%%%%%%%%%%%%%%%%%%%%%%%%%%%%%%%%%%%%%%%%%%
% SUMMARY
%%%%%%%%%%%%%%%%%%%%%%%%%%%%%%%%%%%%%%%%%%%%%%%%%%%%%%%%%%%%%%%%%%%%%%%%%%%%%%%%

\section{Summary}

The dependency ecosystem of the \ProjectName\ consists of third-party libraries that collectively provide the functionality required by the application. The generated Software Bill of Materials (SBOM) offers complete visibility into these dependencies, while the accompanying Grype and ZeroWatch Agent vulnerability assessments evaluate their security posture by identifying known vulnerabilities and available fixes.

Together, dependency analysis, SBOM generation, and vulnerability assessment strengthen software supply chain security, improve governance and compliance, and support proactive vulnerability management throughout the software development lifecycle.